% Standalone problem document, generated from the adjacent README.md.
% Compile with XeLaTeX. Mathematical content is maintained in Markdown.
\documentclass[11pt,a4paper]{article}
\usepackage[fontsize=10.5pt]{scrextend}
\usepackage[margin=22mm,headheight=15pt,headsep=9mm,footskip=11mm]{geometry}
\usepackage{fontspec}
\setmainfont{texgyrepagella}[Extension=.otf,UprightFont=*-regular,BoldFont=*-bold,ItalicFont=*-italic,BoldItalicFont=*-bolditalic]
\setsansfont{texgyreheros}[Extension=.otf,UprightFont=*-regular,BoldFont=*-bold,ItalicFont=*-italic,BoldItalicFont=*-bolditalic]
\setmonofont{texgyrecursor}[Extension=.otf,UprightFont=*-regular,BoldFont=*-bold,ItalicFont=*-italic,BoldItalicFont=*-bolditalic,Scale=MatchLowercase]
\usepackage{amsmath,amssymb}
\usepackage{unicode-math}
\setmathfont{texgyrepagella-math.otf}
\usepackage{xcolor,microtype,enumitem,needspace,fancyhdr}
\usepackage{bookmark}
\definecolor{ink}{HTML}{17344B}
\definecolor{accent}{HTML}{197D80}
\definecolor{muted}{HTML}{596773}
\hypersetup{colorlinks=true,linkcolor=accent,urlcolor=accent,pdftitle={IE-10 - Conditioning
of a random Krylov compression of a cyclic
shift},pdfauthor={Open Problems in Numerical Linear Algebra}}
\urlstyle{same}
\setlength{\parindent}{0pt}
\setlength{\parskip}{5pt plus 1pt minus 1pt}
\linespread{1.04}
\setlength{\emergencystretch}{3em}
\widowpenalty=10000
\clubpenalty=10000
\displaywidowpenalty=10000
\allowdisplaybreaks[1]
\setlist{nosep,leftmargin=1.5em}
\providecommand{\tightlist}{\setlength{\itemsep}{0pt}\setlength{\parskip}{0pt}}
\providecommand{\pandocbounded}[1]{#1}
\pagestyle{fancy}
\fancyhf{}
\fancyhead[L]{\sffamily\footnotesize\color{muted}OPEN PROBLEMS IN NUMERICAL LINEAR ALGEBRA}
\fancyhead[R]{\sffamily\bfseries\footnotesize\color{accent}IE-10}
\fancyfoot[L]{\parbox[b]{0.9\textwidth}{\sffamily\fontsize{8}{10}\selectfont\color{muted}Editorial ratings; a bounded search cannot certify that a problem remains unsolved.\\Literature check: 2026-09-10}}
\fancyfoot[R]{\sffamily\footnotesize\color{muted}\thepage}
\renewcommand{\headrulewidth}{0.3pt}
\renewcommand{\headrule}{\hbox to\headwidth{\color{accent}\leaders\hrule height \headrulewidth\hfill}}
\makeatletter
\renewcommand{\section}{\Needspace{5\baselineskip}\@startsection{section}{1}{0pt}{12pt plus 2pt minus 2pt}{4pt}{\sffamily\bfseries\large\color{ink}}}
\renewcommand{\subsection}{\Needspace{4\baselineskip}\@startsection{subsection}{2}{0pt}{9pt plus 1pt minus 1pt}{3pt}{\sffamily\bfseries\normalsize\color{ink}}}
\makeatother
\setcounter{secnumdepth}{0}
\begin{document}
{\sffamily\small\color{muted}Eigenvalues and inverse problems\par}
\vspace{3pt}
{\raggedright\hyphenpenalty=10000\exhyphenpenalty=10000\sffamily\bfseries\fontsize{21}{24}\selectfont\color{ink}Conditioning
of a random Krylov compression of a cyclic shift\par}
\vspace{7pt}
{\sffamily\small\textbf{Difficulty:} challenging\quad\textbf{Importance:} interesting
to specialist\par}
{\sffamily\small\bfseries\color{accent}Status: Open\par}
\vspace{4pt}
{\color{accent}\hrule height 0.8pt}
\vspace{6pt}
\textbf{Rating rationale:} Challenging reflects a least-singular-value
bound for a dependent random Krylov matrix; specialist impact concerns a
structured compression mechanism in eigenvalue algorithms.

\subsection{Problem statement}\label{problem-statement}

Let \(C_n\in\mathbb C^{n\times n}\) be the cyclic shift,
\(C_ne_j=e_{j+1}\) for \(j<n\) and \(C_ne_n=e_1\). Draw \(b\) uniformly
from the complex unit sphere. For \(2\leq k<n\), let \(Q\) have
orthonormal columns spanning

\[
\mathcal K_k(C_n,b)=\operatorname{span}\{b,C_nb,\ldots,C_n^{k-1}b\},
\qquad H=Q^*C_nQ.
\]

This subspace has dimension \(k\) almost surely. Define
\(\kappa_V(H)=\inf_{H=VDV^{-1},\ D\text{ diagonal}}\|V\|_2\|V^{-1}\|_2\),
and set it to \(+\infty\) if \(H\) is not diagonalizable. Do universal
constants \(C,c>0\) exist such that, for every \(n\geq3\) and
\(2\leq k<n\), with the same constants,

\[
\Pr\{\kappa_V(H)\leq Cn^c\}\geq0.99?
\]

Changing the orthonormal basis of the same Krylov space does not affect
this question. Randomness belongs to the starting vector; replacing
\(Q\) by an independently Haar-distributed subspace would change the
model.

\subsection{Reference}\label{reference}

Amsel et al.,
\href{https://arxiv.org/html/2602.05394v3\#S3.SS3}{\emph{Linear Systems
and Eigenvalue Problems}}, Problem 3.5, which explicitly singles out the
circulant shift. This entry fixes ``high probability'' to a uniform 0.99
success target.

\subsection{Earlier status check ---
2026-09-08}\label{earlier-status-check-2026-09-08}

Searches for \textit{Krylov circulant shift condition 2026} and
\textit{random Krylov compression eigenvector condition number}
found no solution. The deterministic starting vector \(e_1\) gives a
Jordan compression, so the unrandomized statement would be false; that
exceptional example does not settle the probabilistic problem. The
workshop report's 20 August update reports
\href{https://yangpliu.github.io/repository.html}{Peng's obstruction}
for arbitrary real diagonalizable inputs with ill-conditioned
eigenvectors and real Gaussian starts. It explicitly leaves the
normal-input case open, which includes the cyclic shift here.

\subsection{Audit update --- 2026-09-10}\label{audit-update-2026-09-10}

Rechecked \href{https://arxiv.org/html/2602.05394v3}{workshop version
3}, Problem 3.5 and its August update. The arbitrary-diagonalizable
counterexample does not cover this normal cyclic shift. Random
Krylov/circulant conditioning searches found no bound resolving the
displayed structured case.
\end{document}
